\documentclass[UTF8,11pt]{ctexart}
\usepackage[a4paper,margin=24mm]{geometry}
\usepackage{amsmath,amssymb,amsthm,mathtools,mathrsfs}
\usepackage[all]{xy}
\usepackage{xurl,hyperref,enumitem}
\usepackage{tikz}
\usetikzlibrary{arrows.meta,cd,decorations.markings,calc,patterns}
\hypersetup{hidelinks,breaklinks=true}
\setlength{\emergencystretch}{3em}
\allowdisplaybreaks
\newtheorem*{theorem}{Theorem}
\newtheorem*{thm}{Theorem}
\newtheorem*{lemma}{Lemma}
\newtheorem*{proposition}{Proposition}
\newtheorem*{prop}{Proposition}
\newtheorem*{corollary}{Corollary}
\newtheorem*{cor}{Corollary}
\newtheorem*{claim}{Claim}
\theoremstyle{definition}
\newtheorem*{definition}{Definition}
\newtheorem*{defn}{Definition}
\newtheorem*{remark}{Remark}
\newtheorem*{example}{Example}
\newcommand{\R}{\mathbb R}
\newcommand{\C}{\mathbb C}
\newcommand{\Z}{\mathbb Z}
\newcommand{\Q}{\mathbb Q}
\newcommand{\N}{\mathbb N}
\newcommand{\F}{\mathbb F}
\newcommand{\D}{\mathbb D}
\newcommand{\bB}{\mathbb B}
\newcommand{\bC}{\mathbb C}
\newcommand{\bR}{\mathbb R}
\newcommand{\bZ}{\mathbb Z}
\newcommand{\bN}{\mathbb N}
\newcommand{\bQ}{\mathbb Q}
\newcommand{\bD}{\mathbb D}
\newcommand{\bF}{\mathbb F}
\newcommand{\CP}{\mathbb{CP}}
\newcommand{\RP}{\mathbb{RP}}
\newcommand{\sO}{\mathscr O}
\newcommand{\sI}{\mathscr I}
\newcommand{\sF}{\mathscr F}
\newcommand{\sG}{\mathscr G}
\newcommand{\sH}{\mathscr H}
\newcommand{\sR}{\mathscr R}
\newcommand{\sM}{\mathscr M}
\newcommand{\sV}{\mathscr V}
\newcommand{\sU}{\mathscr U}
\DeclareMathOperator{\Hom}{Hom}
\DeclareMathOperator{\Ext}{Ext}
\DeclareMathOperator{\Tor}{Tor}
\DeclareMathOperator{\Spec}{Spec}
\DeclareMathOperator{\Proj}{Proj}
\DeclareMathOperator{\supp}{supp}
\DeclareMathOperator{\im}{im}
\DeclareMathOperator{\id}{id}
\DeclareMathOperator{\Tr}{Tr}
\DeclareMathOperator{\tr}{tr}
\DeclareMathOperator{\rank}{rank}
\DeclareMathOperator{\ord}{ord}
\DeclareMathOperator{\dist}{dist}
\DeclareMathOperator{\End}{End}
\DeclareMathOperator{\Aut}{Aut}
\DeclareMathOperator{\Gal}{Gal}
\DeclareMathOperator{\vol}{vol}
\DeclareMathOperator{\Cl}{Cl}
\DeclareMathOperator{\coker}{coker}
\DeclareMathOperator{\codim}{codim}
\DeclareMathOperator{\divergence}{div}
\DeclareMathOperator{\Ric}{Ric}
\DeclareMathOperator{\grad}{grad}
\DeclareMathOperator{\Hess}{Hess}
\newcommand{\abs}[1]{\left\lvert#1\right\rvert}
\newcommand{\sabs}[1]{\lvert#1\rvert}
\newcommand{\babs}[1]{\bigl\lvert#1\bigr\rvert}
\newcommand{\Babs}[1]{\Bigl\lvert#1\Bigr\rvert}
\newcommand{\bbabs}[1]{\biggl\lvert#1\biggr\rvert}
\newcommand{\BBabs}[1]{\Biggl\lvert#1\Biggr\rvert}
\newcommand{\norm}[1]{\left\lVert#1\right\rVert}
\newcommand{\snorm}[1]{\lVert#1\rVert}
\newcommand{\bnorm}[1]{\bigl\lVert#1\bigr\rVert}
\newcommand{\Bnorm}[1]{\Bigl\lVert#1\Bigr\rVert}
\newcommand{\bbnorm}[1]{\biggl\lVert#1\biggr\rVert}
\newcommand{\BBnorm}[1]{\Biggl\lVert#1\Biggr\rVert}
\newcommand{\widebar}[1]{\overline{#1}}
\newcommand{\nicefrac}[2]{\frac{#1}{#2}}
\newcommand{\linnprod}[2]{\langle#1,#2\rangle}
\newcommand{\myindex}[1]{#1}
\newcommand{\glsadd}[1]{}
\newcommand{\avoidbreak}{}
\newcommand{\sourceref}[1]{\textnormal{[原文：\nolinkurl{#1}]}}
\newcommand{\thmref}[1]{\sourceref{#1}}
\newcommand{\lemmaref}[1]{\sourceref{#1}}
\newcommand{\propref}[1]{\sourceref{#1}}
\newcommand{\corref}[1]{\sourceref{#1}}
\newcommand{\defnref}[1]{\sourceref{#1}}
\newcommand{\exampleref}[1]{\sourceref{#1}}
\newcommand{\exerciseref}[1]{\sourceref{#1}}
\newcommand{\figureref}[1]{\sourceref{#1}}
\newcommand{\sectionref}[1]{\sourceref{#1}}
\newcommand{\appendixref}[1]{\sourceref{#1}}
\newcommand{\stacksref}[2]{\href{https://stacks.math.columbia.edu/tag/#1}{#2 [#1]}}


\newcommand{\E}{\mathbb E}
\newcommand{\Prob}{\mathbb P}
\newcommand{\Reff}{\mathcal R_{\mathrm{eff}}}
\newcommand{\eps}{\varepsilon}
\DeclareMathOperator{\diam}{diam}
\DeclareMathOperator{\Var}{Var}
\DeclareMathOperator{\Crit}{Crit}
\newcommand{\dd}{\,\mathrm d}


\begin{document}
\section*{087\quad 紧群元素的Peter--Weyl有限维酉表示分离}
\subsection*{来源与整理范围}
Terence Tao，254A, Notes 3: Haar measure and the Peter-Weyl theorem，2011-09-27；Section 2，Theorems 6--7 及卷积算子说明。
\par\url{https://terrytao.wordpress.com/2011/09/27/254a-notes-3-haar-measure-and-the-peter-weyl-theorem/}
\par\textbf{恢复方式（整理者说明）：}依据人类原文重新整理以补齐缺失题号；并非旧题证文件的逐字节恢复；2026-09-28。
\subsection*{作者命题与人类证明}

\begin{theorem}[Finite-dimensional separation; Tao's Theorem 7]
Let $G$ be a compact Hausdorff group with Haar probability measure $\mu$.
For every $y\ne e$ there is a finite-dimensional left-invariant subspace
$V\subseteq L^2(G,\mu)$ on which the left translation
$\tau(y)f(x)=f(y^{-1}x)$ is not the identity. In particular, continuous
finite-dimensional unitary representations separate the elements of $G$.
\end{theorem}
\paragraph{Convolution operators.}
The regular representation $\tau:G\to U(L^2(G))$ is strongly continuous.
For $g\in L^2(G)$ let
\[
 T_gf=f*g,\qquad (f*g)(x)=\int_G f(z)g(z^{-1}x)\,d\mu(z).
\]
The convolution belongs to $C(G)$. If
\begin{equation}\label{selfadjoint-kernel}
 g(x^{-1})=\overline{g(x)},
\end{equation}
then $T_g$ is self-adjoint. It commutes with all left translations.
Its integral kernel $K(z,x)=g(z^{-1}x)$ belongs to $L^2(G\times G)$;
indeed $\|K\|_2=\|g\|_2$ by invariance and $\mu(G)=1$.
Thus $T_g$ is compact by the compactness theorem for square-integrable
integral kernels. Each nonzero eigenspace is left-invariant.

\begin{lemma}[Compact self-adjoint spectral theorem; Tao's Theorem 6]
A compact self-adjoint operator $T$ on a complex Hilbert space $H$ has
an orthogonal decomposition
\[
 H=V_0\mathbin{\widehat\oplus}\mathop{\widehat\bigoplus}_{n}V_{\lambda_n},
\]
where $V_0=\ker T$, all $\lambda_n$ are nonzero real eigenvalues, and their
eigenspaces are finite-dimensional. The distinct nonzero eigenvalues are
at most countable and, if infinite in number, tend to zero.
\end{lemma}
\begin{proof}
Self-adjointness makes eigenspaces for distinct eigenvalues orthogonal and
forces eigenvalues to be real. For $r>0$, choose an orthonormal basis of
$\bigoplus_{|\lambda|>r}V_\lambda$. The images under $T$ of distinct basis
vectors are orthogonal and have norms at least $r$. This basis cannot be
infinite, since its image would have no convergent subsequence, contrary
to compactness. Consequently the sum is finite-dimensional for every
$r>0$. This proves finite-dimensionality of every nonzero eigenspace and
the asserted countability and convergence of the eigenvalues.

Let $W$ be the orthogonal complement of the closed span of all eigenspaces,
including the kernel. It is invariant under $T$, and $T|_W$ is self-adjoint
and has no eigenvectors. Suppose $W\ne0$. Then $\|T|_W\|>0$ because the
restriction has trivial kernel. For a bounded self-adjoint operator,
\[
 \|T|_W\|=\sup_{x\in W,\,\|x\|\leq1}|\langle Tx,x\rangle|.
\]
Choose $x_n\in W$, $\|x_n\|\leq1$, such that
$\langle Tx_n,x_n\rangle\to\lambda$, where
$\lambda=\|T|_W\|$ or $\lambda=-\|T|_W\|$. Then
\[
\begin{split}
 0\leq\|Tx_n-\lambda x_n\|^2
 &=\|Tx_n\|^2+\lambda^2\|x_n\|^2
       -2\lambda\langle Tx_n,x_n\rangle\\
 &\leq 2\lambda^2-2\lambda\langle Tx_n,x_n\rangle\longrightarrow0.
\end{split}
\]
Applying $T$ gives $T(Tx_n)-\lambda Tx_n\to0$. By compactness, pass to a
subsequence with $Tx_n\to v$. Then $Tv=\lambda v$. Since $T|_W$ has no
eigenvectors, $v=0$. But this implies
$\langle Tx_n,x_n\rangle\to0$, a contradiction to $\lambda\ne0$.
Thus $W=0$.
\end{proof}

\begin{proof}[Proof of finite-dimensional separation]
Suppose, for a contradiction, that $\tau(y)$ is the identity on every
finite-dimensional invariant subspace. For each $g$ satisfying
\eqref{selfadjoint-kernel}, the nonzero eigenspaces of $T_g$ are such
subspaces. The operator $\tau(y)-1$ annihilates them. It also preserves
$\ker T_g$, since it commutes with $T_g$. The spectral decomposition
therefore gives
\[
 T_g(\tau(y)-1)f=0\qquad(f\in L^2(G)).
\]
Equivalently,
\[
 \tau(y)(f*g)=f*g
\]
for every $f,g\in L^2(G)$ with $g$ satisfying \eqref{selfadjoint-kernel}.
Choose an open symmetric identity neighbourhood $U$ so small that
$y\notin U^2$, and put $f=g=\mathbf1_U$. These functions are square
integrable, $g$ satisfies \eqref{selfadjoint-kernel}, and
\[
 (f*g)(e)=\mu(U)>0,\qquad (f*g)(y)=0.
\]
Since $f*g$ is continuous, the asserted translation equality holds
everywhere, and at $y$ would imply $(f*g)(e)=(f*g)(y)$, a contradiction.
Restriction of the regular representation to the resulting
finite-dimensional invariant subspace is a continuous unitary
representation with nontrivial value at $y$.
\end{proof}

\subsection*{整理说明与先修范围（非作者正文）}
明确收录作者称为 Baby Peter--Weyl 的元素分离定理，不以它冒充带全部重数的完整 Peter--Weyl 分解。紧自伴算子的谱分解证明全部随题收录；不另计该谱定理。原文范数平方展开漏写的 $\lambda$ 已在交叉项补齐；原文混用的 $\rho,\tau$ 统一为左平移 $\tau$。其余是网页公式到独立 TeX 的编辑转录。

先修为 Haar 概率测度及其双不变性、$L^2$ 卷积的连续性、平方可积核算子的紧致性、自伴算子的数值半径公式。难点是选择与群作用交换的紧卷积算子，从其谱空间提取有限维表示，再构造局部支持函数排除某个群元素在所有有限维空间上都不可见。
\subsection*{可修改的简短标注（非作者正文）}
$c$：紧拓扑群的有限维表示分离

$a$：正则表示；Haar测度；卷积核；紧自伴谱分解；有限维不变子空间；局部支持函数

\texttt{cross\_theory\_status = yes}。把拓扑群元素的区别转成平移算子的非平凡性，再通过紧卷积的谱空间获得有限维表示；见最终卷积反证。

全部标注为非穷尽、可修改初稿。
\subsection*{来源许可与编辑归属}
作者 About 页 Copyright etc. 允许带来源引用复用合理篇幅（例如单篇文章）；本题为其指定小节的编辑转录。许可入口：https://terrytao.wordpress.com/about/ 。
ChatGPT 加入题号、中文来源说明、整理范围和初稿标注；编辑说明与人类证明分列。
\end{document}
